%==============================================================================
% FICHAMENTO CRÍTICO — Automatic Jawbone Structure Segmentation on Dental CBCT
% Images via Deep Learning
% Conforme NBR 6023 (referências) e NBR 10520 (citações)
%==============================================================================
\section{Fichamento: \citeonline{tian2024jawbone}}

% --- 1. IDENTIFICAÇÃO ---
\subsection{Identificação}
\begin{tabular}{ll}
\textbf{Título:} & Automatic jawbone structure segmentation on dental CBCT
images via deep learning \\
\textbf{Autor(es):} & Tian, Yuan; Jin, Hao; Wang, Mingzheng; Zhang, Zhejia;
Wang, Ge; Kou, Dazhi \\
\textbf{Periódico:} & Clinical Oral Investigations \\
\textbf{Ano:} & 2024 \\
\textbf{Volume:} & 28 \\
\textbf{Número:} & 12 \\
\textbf{Páginas:} & 663 \\
\textbf{DOI:} & \url{https://doi.org/10.1007/s00784-024-06061-y} \\
\textbf{Qualis:} & A2 (Odontologia) \\
\textbf{Tipo:} & Artigo original \\
\end{tabular}

% --- 2. OBJETIVO ---
\subsection{Objetivo}
Desenvolver e avaliar um sistema de segmentação automática baseado em deep
learning em dois estágios para segmentar o osso cortical mandibular, osso
medular mandibular, osso cortical maxilar e osso medular maxilar em imagens
de tomografia computadorizada de feixe cônico (CBCT), avaliando o impacto de
anomalias de qualidade e comparando com segmentação manual refinada.

% --- 3. METODOLOGIA ---
\subsection{Metodologia}
\textbf{Desenho do estudo:} Estudo experimental de desenvolvimento e validação
de sistema de segmentação baseado em DL. \\
\textbf{Amostra:} 155 exames de CBCT adquiridos com diferentes parâmetros
(campo de visão, resolução, fabricante). \\
\textbf{Métricas principais:} Dice Similarity Coefficient (DSC) e Average
Symmetric Surface Distance (ASSD). \\
\textbf{Validação:} Comparação com ground truth manual; análise de subgrupo por
anomalias de qualidade (artefatos de movimento, ruído, hipodensidades);
comparação entre segmentação automática (AS) e segmentação manual refinada (MRS).

% --- 4. RESULTADOS PRINCIPAIS ---
\subsection{Resultados Principais}
\begin{itemize}
  \item DSC médio: mandíbula cortical = 93,69\%; mandíbula medular = 96,83\%;
        maxila cortical = 86,14\%; maxila medular = 95,57\%.
  \item ASSD médio: mandíbula cortical = 0,13 mm; mandíbula medular = 0,16 mm;
        maxila cortical = 0,29 mm; maxila medular = 0,41 mm.
  \item A segmentação da maxila cortical apresentou o menor DSC (86,14\%) e o
        maior ASSD (0,29 mm), refletindo a complexidade anatômica e a fina
        espessura cortical na maxila.
  \item A comparação AS $\times$ MRS mostrou DSC $>$ 98,8\% e ASSD $<$ 0,1 mm,
        indicando alta concordância entre a segmentação automática e a refinada.
\end{itemize}

% --- 5. LIMITAÇÕES ---
\subsection{Limitações}
\begin{itemize}
  \item Anomalias de qualidade (movimento, ruído, hipodensidades) reduziram
        significativamente o desempenho, especialmente na maxila cortical.
  \item A amostra de 155 exames é relativamente modesta considerando a
        diversidade de parâmetros de aquisição e variações anatômicas.
  \item O estudo não incluiu segmentação de estruturas clinicamente relevantes
        como seios maxilares, canal mandibular ou dentes.
\end{itemize}

% --- 6. CONTRIBUIÇÃO PARA O LIVRO ---
\subsection{Contribuição para o Livro}
Este artigo é a principal referência do Capítulo 19, que aborda o uso de U-Net
para segmentação óssea em CBCT. A arquitetura em dois estágios (detecção da
região de interesse seguida de segmentação refinada) é apresentada no livro
como uma variação da arquitetura U-Net clássica adaptada para imagens
volumétricas. A análise do impacto de anomalias de qualidade no desempenho
fornece subsídios para a discussão sobre robustez e pré-processamento de
imagens CBCT. O DSC de 96,83\% para osso medular mandibular é utilizado como
benchmark no livro para comparação com outras arquiteturas de segmentação.

% --- 7. ANÁLISE CRÍTICA ---
\subsection{Análise Crítica}
\begin{itemize}
  \item \textbf{Validade interna:} Alta -- Uso de 155 exames com parâmetros
        variados, validação com ground truth manual e análise de subgrupos.
  \item \textbf{Validade externa:} Média -- Dados de múltiplos parâmetros, mas
        não há validação multicêntrica ou em populações diversas.
  \item \textbf{Reprodutibilidade:} Média -- A arquitetura em dois estágios é
        descrita conceitualmente, mas os autores não disponibilizaram código
        ou pesos pré-treinados publicamente.
  \item \textbf{Relevância clínica:} Alta -- Segmentação óssea automática é
        pré-requisito para planejamento de implantes, cirurgia ortognática e
        geração de gêmeos digitais.
  \item \textbf{Conflito de interesses:} Não -- Os autores declararam ausência de
        conflitos.
  \item \textbf{Viés identificado:} Viés de espectro -- A amostra pode não
        representar toda a gama de variações anatômicas e patológicas
        encontradas na prática clínica.
\end{itemize}

% --- 8. CITAÇÕES RELEVANTES ---
\subsection{Citações Relevantes}
\begin{citacao}
``The system achieved promising segmentation performance, with average DSC
values of 93.69\%, 96.83\%, 86.14\% and 95.57\% and average ASSD values of
0.13 mm, 0.16 mm, 0.29 mm and 0.41 mm for the mandibular cortical bone,
mandibular cancellous bone, maxillary cortical bone and maxillary cancellous
bone, respectively.'' (p. 663).
\end{citacao}

% --- 9. MAPEAMENTO SDD ---
\subsection{Mapeamento SDD}
\textbf{Problema:} Segmentação manual de estruturas ósseas maxilomandibulares
em CBCT é demorada, subjetiva e impraticável em larga escala clínica. \\
\textbf{Entrada:} Imagens CBCT (formato DICOM) com parâmetros de aquisição
variáveis (FOV, resolução, fabricante), 155 exames. \\
\textbf{Saída:} Mapas de segmentação com 4 classes: osso cortical mandibular,
osso medular mandibular, osso cortical maxilar, osso medular maxilar. \\
\textbf{Critérios:} DSC $\geq$ 85\% para todas as classes; ASSD $\leq$ 0,5 mm;
robustez a anomalias de qualidade (movimento, ruído, hipodensidades). \\
\textbf{Confiança:} 0,78 -- Estudo tecnicamente sólido, mas com amostra
moderada, sem código aberto e sem validação externa multicêntrica.

% --- 10. REFERÊNCIA ABNT ---
\subsection{Referência ABNT}
\begin{verbatim}
TIAN, Yuan et al. Automatic jawbone structure segmentation on dental CBCT
images via deep learning. Clinical Oral Investigations, Berlin, v. 28, n. 12,
p. 663, 2024. DOI: 10.1007/s00784-024-06061-y.
\end{verbatim}
